\section{결론}
표~\ref{tab:goal}에 목표 달성 현황을 정리하였다.

\begin{table}[!t]
\centering\footnotesize\setlength{\tabcolsep}{3pt}
\caption{목표 달성 현황}\label{tab:goal}
\begin{tabularx}{\linewidth}{@{}>{\raggedright\arraybackslash}p{21mm}p{7mm}Y@{}}
\toprule
목표 & 결과 & 근거 \\
\midrule
실패 발견율 30\% 이상 & \textbf{달성} & 로봇팔 100\% (6/6); 휴머노이드 33.3\% (6/18) \\
Random 테스트 대비 20\% 이상 향상 & \textbf{달성} & 로봇팔 +500\% (16.7\%$\to$100\%), 휴머노이드 +200\% (11.1\%$\to$33.3\%) \\
실패 유형 6종 중 4종 & \textbf{달성} & 충돌, 장애물 간섭, 목표 위치 점유, 작업자 안전 위험 + 도달 불가 부분 관측 \\
3D 장면 구조화·후보 생성 & 구현 & 10개·21개 장면 변수, 장면 그래프·시뮬레이션 모델 생성, 실행 전 검사 \\
행동 기반 취약 조건 탐색 & 구현 & 관측$\to$가설$\to$제안$\to$FAIL 연결 확인 \\
회귀 테스트 자산화 & 구현 & 실행 기록 저장, 재실행·비교 실행기, 로봇 LLM 교체 후 저장 사례 결과 유지 확인 \\
\bottomrule
\end{tabularx}
\end{table}


VLA 기반 로봇팔에서 AFS는 같은 시도 횟수에서 Random 테스트보다 6배, 휴머노이드에서는 3세트
모두에서 일관되게 Random 테스트보다 3배 많은 실패 사례를 찾았다. 또한 장애물 높이와 접시 위치에 대한
PASS/FAIL 경계를 정량적으로 제시하여, 배포 전 검증에서 어떤 조건을 우선 점검해야 하는지 보여 준다.
LLM의 실패 가설은 관측된 실패 메커니즘(장애물 차단, 목표 점유)을 겨냥할 때 실제 결과와 잘 맞았으나,
아직 관측되지 않은 유형(충돌, 이동 정체)을 겨냥할 때는 빗나갔다.

현재 실험은 실험 세트 수와 반복 수가 작고, AFS가 동일 환경을 반복 제안하는 경향이 있다.
후속 과제로 (1) 더 많은 실험 세트에서의 독립 평가, (2) 새로운 환경 탐색을 보상하는 다양성 목표,
(3) 유사 물체·카메라 가림 조건을 통한 인식 오류 탐색, (4) 반복 수를 늘린 회귀 재실행과 판단 횟수 등 보조 지표 변화 검출을 진행한다.
